\paragraph{\S9-fragment-statement.}
\paragraph{Motivating conjecture, restated verbatim.}
Under Assumptions~\ref{ass:causal}--\ref{ass:effective}, the estimator $\hat\theta_n^{\mathrm{locDR}}$ satisfies
\[
  \left|\hat\theta_n^{\mathrm{locDR}}-\tau_P(x_0)\right|
  =O_P\!\left(b_n+(n p_n)^{-1/2}+\sqrt{r_{g,n}r_{m,n}}\right)
\]
uniformly over $\mathcal{P}_n^{\mathrm{SA\text{-}CATE}}$.

\begin{theorem}[Conjecture 1: upper bound for the localized DR learner]\label{thm:upper}
Under Assumptions~\ref{ass:causal}--\ref{ass:effective}, with the approved clipped-score evaluation convention
\[
  \hat g^c_{a,-k}=\Pi_{[-M,M]}\hat g_{a,-k},\qquad
  \hat e^c_{-k}=\Pi_{[\epsilon/2,\,1-\epsilon/2]}\hat e_{-k},
\]
where $\Pi_{[a,b]}(u)=\min\{b,\max\{a,u\}\}$ and the local nuisance-rate bounds in Assumption~\ref{ass:nuisance} are understood for these clipped score inputs, the explicit cross-fit estimator
\[
  \hat\theta_n^{\mathrm{locDR}}
  =
  \frac{1}{n}\sum_{k=1}^K\sum_{i\in I_k}
  \psi_{h_n}\!\left(W_i;(\hat g^c_{0,-k},\hat g^c_{1,-k},\hat e^c_{-k})\right)
\]
satisfies $\left|\hat\theta_n^{\mathrm{locDR}}-\tau_P(x_0)\right|=O_P\!\left(b_n+(n p_n)^{-1/2}+\sqrt{r_{g,n}r_{m,n}}\right)$ uniformly over $\mathcal{P}_n^{\mathrm{SA\text{-}CATE}}$.
\end{theorem}

\paragraph{\S13-refutation-attempt.}
\paragraph{Primitive attempted.}
The refutation search used the influence-function/variance primitive.  The searched witnesses were the smallest finite submodels that could break the stated rate: an oracle-nuisance submodel, a clipped estimated-nuisance submodel satisfying the local $L_2$ rate envelopes, and a boundary submodel in which unclipped estimated propensities approach zero on a shrinking local set.

\paragraph{Oracle-nuisance witness.}
Set $\hat\eta_{-k}=\eta_P$ for every fold.  Then Theorem~\ref{thm:localized-aipw} gives $E_P[\psi_{h_n}(W;\eta_P)]=\theta_{h_n}(P)$ and
\[
  \hat\theta_n^{\mathrm{locDR}}-\theta_{h_n}(P)
  =
  \frac1n\sum_{i=1}^n\{\psi_{h_n}(W_i;\eta_P)-E_P\psi_{h_n}(W;\eta_P)\}.
\]
Assumptions~\ref{ass:bounded} and~\ref{ass:effective} imply
$\operatorname{Var}_P\{\psi_{h_n}(W;\eta_P)\}\lesssim p_n^{-1}$, so this witness only produces the allowed stochastic scale $(n p_n)^{-1/2}$.

\paragraph{Clipped nuisance witness.}
Now allow estimated nuisances satisfying Assumption~\ref{ass:nuisance}, but evaluate the score with the clipped inputs displayed in Theorem~\ref{thm:upper}.  The product-drift half of Theorem~\ref{thm:localized-aipw} bounds the foldwise conditional mean drift by
\[
  C_\epsilon R_{g,n,k}R_{m,n,k}
  =
  O_P\!\left(\sqrt{r_{g,n}r_{m,n}}\right).
\]
The clipped denominators and bounded outcomes keep the conditional second moment at order $p_n^{-1}$, so this candidate also matches, rather than violates, the conjectured rate.

\paragraph{Boundary search.}
The only boundary pattern that can exceed the rate is an untruncated estimated propensity that is nearly zero on a set of tiny local mass while its local $L_2$ distance from $e$ remains small.  That construction is excluded by the approved clipped-score evaluation convention, which is now part of the estimator convention for this correction round.  Hence no counterexample remains under the stated and approved regime, and the positive proof proceeds.

\paragraph{\S10 Interpretation.}
The theorem is the achievability half of the structure-agnostic pointwise CATE rate.  The estimator is an explicit cross-fitted localized AIPW average, so its error separates into three channels: localization bias $b_n$, oracle local empirical noise $(n p_n)^{-1/2}$, and the second-order nuisance product $\sqrt{r_{g,n}r_{m,n}}$.  This is the normalized point-local analogue of the doubly robust pseudo-outcome decomposition in \cite{Kennedy2023DRCATE}: cross-fitting removes own-observation bias, orthogonality eliminates first-order nuisance drift, and the localizer inflates variance by $p_n^{-1}$ because only $n p_n$ observations effectively contribute near $x_0$.

\paragraph{\S11 Proof.}
\begin{proof}[Proof of Theorem~\ref{thm:upper}]
Write $\ell_n=\ell_{h_n}$, $\theta_n(P)=\theta_{h_n}(P)$, and
\[
  \hat\eta^c_{-k}=(\hat g^c_{0,-k},\hat g^c_{1,-k},\hat e^c_{-k}).
\]
The projection map onto a closed interval is nonexpansive.  Since Assumption~\ref{ass:causal} gives
$e\in[\epsilon,1-\epsilon]$ and Assumption~\ref{ass:bounded} gives $g_a\in[-M,M]$, clipping cannot
increase the local weighted nuisance errors:
\[
  E_P[\ell_n(X)(\hat e^c_{-k}(X)-e(X))^2]
  \le E_P[\ell_n(X)(\hat e_{-k}(X)-e(X))^2],
\]
and the same inequality holds for $\hat g^c_{a,-k}-g_a$, $a\in\{0,1\}$.  Thus the clipped score inputs
obey the same rate envelopes $r_{g,n}$ and $r_{m,n}$ as the local nuisance bounds in
Assumption~\ref{ass:nuisance}.

Decompose the estimator around the smoothed target:
\[
\begin{aligned}
  \hat\theta_n^{\mathrm{locDR}}-\theta_n(P)
  &=
  \underbrace{
  \frac{1}{n}\sum_{k=1}^K\sum_{i\in I_k}
  \left\{\psi_{h_n}(W_i;\hat\eta^c_{-k})
       -E_P[\psi_{h_n}(W;\hat\eta^c_{-k})\mid \hat\eta^c_{-k}]\right\}}_{A_n} \\
  &\quad+
  \underbrace{
  \frac{1}{n}\sum_{k=1}^K |I_k|
  \left\{E_P[\psi_{h_n}(W;\hat\eta^c_{-k})\mid \hat\eta^c_{-k}]-\theta_n(P)\right\}}_{B_n}.
\end{aligned}
\]
Conditional on the training folds, the observations in $I_k$ are independent of
$\hat\eta^c_{-k}$, so the summands in $A_n$ have conditional mean zero.  The clipped-score convention,
Assumption~\ref{ass:bounded}, and $E_P[\ell_n(X)^2]\lesssim p_n^{-1}$ imply that, for a constant $C$
depending only on $M$ and $\epsilon$,
\[
  E_P\!\left[\psi_{h_n}(W;\hat\eta^c_{-k})^2\mid \hat\eta^c_{-k}\right]
  \le C E_P[\ell_n(X)^2]
  \le C p_n^{-1}.
\]
Therefore
\[
  \operatorname{Var}_P(A_n\mid \hat\eta^c_{-1},\ldots,\hat\eta^c_{-K})
  \le
  \frac{1}{n^2}\sum_{k=1}^K |I_k|\,C p_n^{-1}
  \le \frac{C}{n p_n}.
\]
Chebyshev's inequality yields $A_n=O_P((n p_n)^{-1/2})$ uniformly, because the constants are uniform
over $\mathcal{P}_n^{\mathrm{SA\text{-}CATE}}$ and $n p_n\to\infty$ by
Assumption~\ref{ass:effective}.

For the drift term $B_n$, apply Theorem~\ref{thm:localized-aipw} fold by fold.  Since
$\hat e^c_{-k}\in[\epsilon/2,1-\epsilon/2]$ by construction,
\[
  \left|
  E_P[\psi_{h_n}(W;\hat\eta^c_{-k})\mid \hat\eta^c_{-k}]
  -\theta_n(P)
  \right|
  \le
  C_\epsilon R_{g,n,k}R_{m,n,k},
\]
where
\[
  R_{g,n,k}^2
  =
  E_P[\ell_n(X)\{(\hat g^c_{1,-k}-g_1)^2+(\hat g^c_{0,-k}-g_0)^2\}],
  \qquad
  R_{m,n,k}^2
  =
  E_P[\ell_n(X)(\hat e^c_{-k}-e)^2].
\]
Assumption~\ref{ass:nuisance} and the fixed number of folds imply
$\max_k R_{g,n,k}=O_P(r_{g,n}^{1/2})$ and
$\max_k R_{m,n,k}=O_P(r_{m,n}^{1/2})$.  Since $\sum_k |I_k|/n=1$,
\[
  |B_n|
  \le C_\epsilon \max_{1\le k\le K} R_{g,n,k}R_{m,n,k}
  =
  O_P\!\left(\sqrt{r_{g,n}r_{m,n}}\right)
\]
uniformly over the model class.

Finally, Assumption~\ref{ass:bias} gives
\[
  |\theta_n(P)-\tau_P(x_0)|\le b_n.
\]
Combining the empirical fluctuation, drift, and localization-bias bounds,
\[
  |\hat\theta_n^{\mathrm{locDR}}-\tau_P(x_0)|
  \le |A_n|+|B_n|+|\theta_n(P)-\tau_P(x_0)|
  =
  O_P\!\left(b_n+(n p_n)^{-1/2}+\sqrt{r_{g,n}r_{m,n}}\right).
\]
This proves the theorem.  The argument is the DR-CATE pseudo-outcome decomposition of
\cite{Kennedy2023DRCATE} with the normalized local weight $\ell_n=K_{h_n}/p_n$ making the effective
sample size $n p_n$ explicit.
\end{proof}

\paragraph{\S12 Sanity check.}
Take the global randomized corner case $K_{h_n}(X)\equiv1$, so $p_n=1$ and $\ell_n\equiv1$, and suppose
$\tau_P(x)=\tau_0$ is constant so $b_n=0$.  If the nuisances are known, then
$r_{g,n}=r_{m,n}=0$ and the estimator reduces to the usual AIPW average for the average treatment
effect.  The theorem gives
\[
  |\hat\theta_n^{\mathrm{locDR}}-\tau_0|=O_P(n^{-1/2}),
\]
which is the classical root-$n$ AIPW scale under bounded outcomes and overlap.

\paragraph{\S13 Counterexample or minimality check.}
The clipped-score convention is the minimal variance-control condition for this upper bound.  To see
what fails without it, let the true propensity be $e(X)=1/2$, let $g_0(X)=g_1(X)=0$, and let $Y$ be
bounded with nonzero conditional variance among treated units.  Let $K_{h_n}$ be the indicator of a
local set $A_n$ with $P(X\in A_n)=p_n$, and choose $B_n\subset A_n$ with $P(X\in B_n)=1/n$.  Define a
raw fold-specific propensity estimate by $\hat e_{-k}(x)=\delta_n$ on $B_n$ and
$\hat e_{-k}(x)=1/2$ elsewhere, with $\delta_n\downarrow0$ much faster than $(n p_n)^{-1}$.
Then
\[
  E_P[\ell_n(X)(\hat e_{-k}(X)-e(X))^2]\lesssim \frac{1}{n p_n}\to0,
\]
so the local $L_2$ nuisance condition can still hold at a small rate.  However, on the event that one
evaluation observation falls in $B_n$, is treated, and has a nonzero residual, the single treated
residual term contributes order
\[
  \frac{1}{n}\frac{\ell_n(X)}{\delta_n}
  =
  \frac{1}{n p_n\delta_n},
\]
which can be made arbitrarily larger than $(n p_n)^{-1/2}$.  This boundary example is excluded by the
approved clipping convention and explains why the score-denominator bound is not cosmetic.
